% généré par scripts/06_textes.py
\dosmeta{Année de classement : 389 (directe : nom consulaire lu). Références : CIL VI 41342 = ICUR V 13355 = AE 2009, 145. Rattachement selon l’EDB : \emph{Coem. subdiale ad Catacumbas}. Conservation selon l’EDB : \emph{Coem. subdiale ad Catacumbas}. Support : Sarcophage de marbre. 72 × 212 cm (h. × l.). Inscription gravée. Reproductions signalées par l’EDB : photographie PCAS ×2 (droits réservés ; non reproduites). Fiche EDB : Anita Rocco, première version 2010. Notice du catalogue p.~\pageref{not:306}.}
\txtbloc{Texte transmis (EDCS, état 2022)}{\lnum{1}Maximilla poten[s servare f(o)edera C]⸢hr⸣isti propositumque fidem mors abstulit omnia vitae sed melius vivit vitam C(h)risti qui sperat (a)eternam si animo recolat fidem datam omn[ibus] unam tunc poterit mors ipsa mori cum tempore toto vita perennis eri⸢t⸣ veniam dabit omnibus unam sed meritis [ve]⸢g⸣etans su[per h]anc Lucceia fecisti nam m⸢e⸣mor ipsa tui cum reddis vota sepulcr\{h\}i ostend⸢i⸣s mentem religiosa vita repleta(m) mirae fide(i) et innocentiae optimique propositi Maximilla virgo ancilla dei civis Pannonia qu(a)e vixit annis n(umero) L huic sepulturae condita est a \{L\} Lucceia c(larissima) f(emina) filia Viventi c(larissimae) r(ecordationis) v(iri) ex praefecto pr(a)etorio et urbis aeternae prop⸢te⸣r amicitiam sibi coniun(c)tam qu(a)e diem clusit una cum matre sua pariter Nunita qu(a)e fuit matrona diaconi\{s\} et vixit annis n(umero) LXX die IIII Non(as) Octobr(es) FFll(aviis) Timasio et Promoto vv(viris) cc(larissimis) conss(ulibus)}
\txtbloc{Texte de l’EDB (EDB27045)}{\lnum{1}Maximilla poten[s servare foedera C]rhisti  propositumque fidem mors abstulit omnia vitae sed melius vivit vitam Cristi qui sperat eternam si animo recolat fidem\par
\lnum{2}datam omn[ibus] unam tunc poterit mors ipsa mori cum tempore toto  vita perennis eri⌜t⌝ veniam dabit omnibus unam sed meritis\par
\lnum{3}[ve]getans su[per h]anc Lucceia fecisti nam memor ipsa tui cum reddis vota sepulcrhi ostendis mentem religiosa vita repleta⟨m⟩. Mirae fide⟨i⟩\par
\lnum{4}et innocentiae optimique propositi Maximilla virgo ancilla dei civis Pannonia que\par
\lnum{5}vixit annis n(umero) L huic sepulturae condita est a \{L\} Lucceia c(larissima) f(emina) filia Viventi c̅(larissimae) r̅(ecordationis) v̅(iri) ex praefecto\par
\lnum{6}pretorio et urbis aeternae propter amicitiam sibi coniuntam que diem clusit una cum matre\par
\lnum{7}sua pariter Nunita que fuit matrona diaconis et vixit annis n(umero) LXX die IIII non(as) octobr(es) Ff(laviis duobus)\par
\lnum{8}Timasio et Promoto vv(ris) cc(larissimis duobus) conss(ulibus duobus)}
